\documentclass[10pt,a4paper]{amsart}
\usepackage[margin=19mm]{geometry}
\usepackage{amsmath,amssymb,mathtools}
\usepackage[scr=rsfs,cal=euler]{mathalfa}
\usepackage{xfrac}
\usepackage[unicode,hidelinks,bookmarks=false]{hyperref}
\newcommand{\ie}{i.e.\ }
\newcommand{\cf}{cf.\ }
\newcommand{\resp}{respectively\ }
\newcommand{\Subsection}{\subsection}
\providecommand{\nobreakdash}{\nobreak\hbox{-}}
\setlength{\emergencystretch}{2em}
\multlinegap0pt
\usepackage{mathtools}
\usepackage{amssymb}
\usepackage[shortlabels]{enumitem}
\usepackage{tikz}
\usepackage{xcolor}
\newtheorem{claim}{Claim}
\newtheorem{lemma}{Lemma}
\newcommand{\Gnp}{G_{n, p}^{(3)}}
\renewcommand{\Pr}{\mathbb{P}}
\newcommand{\cF}{\mathcal{F}}
\newcommand{\cFQL}{\cF_Q^{\mathrm{low}}}
\newcommand{\cQ}{\mathcal{Q}}
\newcommand{\ext}{\mathrm{ext}}
\renewcommand{\ge}{\geqslant}
\renewcommand{\le}{\leqslant}
\title{Random Tur\'an Theorem for the Fano Plane}
\author{Ilay Hoshen}
\date{Source: arXiv:2607.28071v1 (CC BY 4.0)}
\begin{document}
\maketitle

\noindent\emph{Source and licence.} Random Tur\'an Theorem for the Fano Plane. Version arXiv:2607.28071v1 (CC BY 4.0), distributed by arXiv under the Creative Commons Attribution 4.0 International licence, http://creativecommons.org/licenses/by/4.0/, as recorded in the arXiv licence metadata for this exact version and shown on the abstract page https://arxiv.org/abs/2607.28071v1. Full licence notice: \texttt{licenses/arxiv-2607.28071-CC-BY-4.0.txt}.\par\smallskip

\noindent\textit{Editorial scope.} Counted unit: the article's lemma lemma:mu-Delta-bounds-low-degree (source0.tex lines 1364--1373). The article's complete original proof of it is delivered with it (source0.tex lines 1374--1480, one \texttt{proof} environment of 107 lines). No mathematics is written, repaired or re-derived: the statement and the proof are the article's own lines, in the article's own notation, and no retained line is substituted or re-flowed.  Retained uncounted as the source-local context the proof needs: lines 1145--1155.  One declared adaptation: exactly 1 ancillary strings inside the retained spans are omitted (image-inclusion commands and, where the source repeats a label, the repeated label definitions) and no other line differs from the source; the figure environments, with their placement, captions and labels, are kept, and the package records the omitted strings in fidelity receipts bound to the affected bodies.  No retained line contains a citation, so the wrapper carries no bibliography and no bibliography entry.  The retained lines contain the article's own diagram code, which the wrapper typesets with the standard tikz packages; no image file is delivered and the package declares no asset dependency.  Editorial formatting only, in addition: an AMS \texttt{amsart} preamble that restates the article's own declarations and macros used by the retained lines, namely the environments \texttt{claim}, \texttt{lemma} on separate counters, together with the standard packages those lines need.  The retained environments print in this document as Claim 1, Lemma 1 in order of appearance, whereas the article prints the same items with the numbering of its own full document.  The complete native source of the article as distributed in the arXiv e-print for this exact version is bundled unchanged as \texttt{sources/206384/source0.tex}.
\begin{claim}\label{claim:strictly-balanced}
    Let $\ell \ge 2$ be an integer and let $H$ be an $\ell$-uniform $\ell$-balanced graph. Suppose that $p \ge C n^{-1/m_\ell(H)}$ for some $C > 0$. Then,
    \[
        n^{v(\tilde{H})} p^{e(\tilde{H})} \ge C^{e(\tilde{H}) - 1} n^{\ell} p,
    \]
    for every $\emptyset \neq \tilde{H} \subseteq H$. Moreover, if $H$ is also stricty $\ell$-balanced, there exists $\lambda > 0$ such that 
    \[
        n^{v(\tilde{H})} p^{e(\tilde{H})} \ge C^{e(\tilde{H}) - 1} n^{\ell + \lambda} p,
    \]
    for every $\emptyset \neq \tilde{H} \subseteq H$ with $\ell < v(\tilde{H}) < v(H)$.
\end{claim}

\begin{lemma}\label{lemma:mu-Delta-bounds-low-degree}
    There exists a positive constant $C_{Low} > 0$ such that, for every sufficiently small $\tilde{\epsilon} > 0$, the following folds for every $p \ge n^{-2/3}$, every $2$-tuple $\Pi = (A_1, A_2)$ of disjoint sets of vertices of size at least $n/4$ each and for every $Q \in \cQ_1 \cup \cQ_2$ with $V(Q) \subseteq A_1$.
    \begin{enumerate}
        \item $\mu_p\left(\cFQL[\ext(\Pi)]\right) \ge \left(\frac{5}{4} - o(1)\right) \cdot e(Q) \cdot \left(\min\left\{|A_1|, |A_2|\right\}\right)^{4} \cdot p^6.$
        \item There exists $\lambda > 0$ such that 
        \begin{align*}
            \frac{\Delta_p\left(\cFQL\right)}{\mu_p\left(\cFQL[\ext(\Pi)]\right)} \le \frac{C_{Low}}{3}\left(\Delta_1(Q) \cdot n^{2} p^{5} + \Delta_2(Q) \cdot n^{3 - \lambda} p^{5} + n^{4 - \lambda} p^{6}\right) \le \frac{C_{Low} \cdot \tilde{\epsilon} \cdot n^4 p^6}{\log n}.
        \end{align*}
    \end{enumerate}
\end{lemma}

\begin{proof}
    For the first item, we have 
    \[
        \mu_p\left(\cFQL[\ext(\Pi)]\right) = \sum_{F' \in \cFQL[\ext(\Pi)]} \Pr\left(F' \subseteq \Gnp\right) = \left|\cFQL[\ext(\Pi)]\right| \cdot p^6,
    \]
    where the last equality follows from the fact that $\cFQL[\ext(\Pi)]$ is a family of subgraphs of $K_n^{(3)} \setminus Q$, each consisting of exactly six edges. To estimate $\left|\cFQL[\ext(\Pi)]\right|$, note that for every $F' \in \cFQL[\ext(\Pi)]$, there exists some $f \in Q$ such that $F' \cup \{f\} \cong F$. There are $e(Q)$ choices for such an edge. Once $f \in Q$ is fixed, we consider the possible choices for the remaining four vertices of $F'$. Observe that $F'$ can take one of the following two forms:
    it either uses one additional vertex from $A_1$ and  three from $A_2$, or it uses four vertices from $A_2$. Furthermore, given an edge of $Q$ and four such additional vertices, there are $3!$ ways to form a copy of $F'$ such that $F' \cup \{f\} \cong F$ using those seven vertices. This is illustrated in Figure \ref{figure:copies of F}.
    
\begin{figure}[h]
    \centering
           
    %\end{center}
    \caption{The figure on the left presents a copy of $F$ lying on an edge $f \in Q$ with four additional vertices in $A_2$. The blue vertex, together with the two vertices in the first blue set in $A_2$, forms an edge, as does the blue vertex with the two vertices in the second blue set. The same structure applies to the green and red vertices with their corresponding sets of size two in $A_2$ of matching colours. Observe that there are $3!$ ways to construct a copy of the Fano plane in such a way.}
    \label{figure:copies of F}
\end{figure}

    % Further, given an edge of $Q$ and another such four vertices, there are $3!$ ways to form a copy of the $F'$ such that $F' \cup e \cong F$ using those seven vertices. Indeed, there are three possible partitions of the other four vertices into two sets of two vertices. In order to form $F$, we should choose, for every vertex in the edge of $Q$, a distinct partition and add the two edges of this vertex and the two sets in the chosen partition. There are $3!$ possibilities to choose a distinct partition for every vertex of the edge of $Q$. 
    
    Once we fix such a copy $F' \in \cFQL[\ext(\Pi)]$, the probability that it is contained in $\Gnp$ is $p^6$.   
    Hence,
    \begin{align*}
        \mu_p\left(\cFQL[\ext(\Pi)]\right) &= (1+o(1)) e(Q) \left[\binom{|A_2|}{4} \cdot 3! \cdot p^6 + \binom{|A_2|}{3} \cdot (|A_1| - 3) \cdot 3! \cdot p^6\right] \\
        &\ge \left(\frac{5}{4} - o(1)\right) \cdot e(Q) \cdot \left(\min\left\{|A_1|, |A_2|\right\}\right)^{4} \cdot p^6.
    \end{align*}    

    Moving to the second item of the lemma, we have
    \begin{align*}
        \Delta_p\left(\cFQL\right) &= \sum_{F' \in \cFQL} \Pr\left(F' \subseteq \Gnp\right) \sum_{\substack{F' \neq F'' \in \cFQL \\ F' \cap F'' \neq \emptyset}} \Pr\left(F'' \subseteq \Gnp \mid F' \subseteq \Gnp\right) \\
        &\le \left|\cFQL\right| \cdot p^{6} \cdot \max_{F' \in \cFQL} \sum_{\substack{F' \neq F'' \in \cFQL \\ F' \cap F'' \neq \emptyset}} p^{6 - e(F' \cap F'')}\\
        &\le \sqrt{C_{Low}} \cdot \mu_p\left(\cFQL[\ext(\Pi)]\right) \cdot \max_{F' \in \cFQL} \sum_{\substack{F' \neq F'' \in \cFQL \\ F' \cap F'' \neq \emptyset}} p^{6 - e(F' \cap F'')}.
    \end{align*}
    where the last inequality is true for large enough $C_{Low}$ since 
    \[
        \left|\cFQL\right| p^{6} = \Theta\left(e(Q) \cdot n^{4} p^{6}\right) = O\left(\mu_p\left(\cFQL[\ext(\Pi)]\right)\right).
    \]
    We will give an upper bound to the above sum for every $F' \in \cFQL$.    

    Fix $F' \in \cFQL$. For every $F'' \in \cFQL$, there exists $f'' \in Q$ such that $F'' \cup \{f''\} \cong F$. We have
    \begin{align}\label{align:sum-of-low-delta}
        \sum_{\substack{F' \neq F'' \in \cFQL \\ F' \cap F'' \neq \emptyset}} p^{6 - e(F' \cap F'')} = 
        \sum_{i=0}^{3} \sum_{\substack{f'' \in Q \\ |f'' \cap V(F')| = i}} \sum_{\substack{F'' \in \cFQL \\ F'' \cup \{f''\} \cong F \\ F' \cap F'' \neq \emptyset}} p^{6 - e(F' \cap F'')}.
    \end{align}
    We analyse the above sum for each case $i \in \{0, 1, 2, 3\}$, where $i$ denotes the number of vertices in the intersection $f'' \cap V(F')$. For the case $i=0$, we observe that no such copies $F''$ exist. Indeed, otherwise $F'$ would contain two vertex-disjoint edges that contradict the fact that every pair of edges in the Fano plane intersects.

    In the case $i = 1$, there are $O(1)$ ways to choose the intersection $\tilde{F} \coloneqq F' \cap F''$. Following the second sum in \eqref{align:sum-of-low-delta}, there are $O(\Delta_1(Q))$ possibilities for the edge $f'' \in Q$. To complete the copy of $F''$, we must select $5 - v(\tilde{F})$ additional vertices and $6 - e(\tilde{F})$ additional edges. Since $F$ is strictly $3$-balanced and $p \ge n^{-1/m_3(F)}$, then, by Claim \ref{claim:strictly-balanced}, we have $n^{v(\tilde{F})} p^{e(\tilde{F})} \ge n^3 p$, implying that $n^{5 - v(\tilde{F})} p^{6 - e(\tilde{F})} \le n^{2} p^{5}$. 
    Hence, the contribution of the $i=1$ case to the sum in \eqref{align:sum-of-low-delta} is
    \[
        O\left(\Delta_1(Q) \cdot n^{2} p^{5}\right) = O\left(\tilde{\epsilon} \cdot \frac{n^{4} p^{6}}{\log n}\right),
    \]
where we used that $\Delta_1(Q) \le \tilde{\epsilon} \frac{n^2 p}{\log n}$ which holds by the definition of $\cQ_1$ and $\cQ_2$.

    In the case $i=2$, there are $O(1)$ possibilities for the intersection $\tilde{F} \coloneqq F' \cap F''$. From the second sum in \eqref{align:sum-of-low-delta}, we have $O(\Delta_2(Q))$ choices for the edge $f'' \in Q$. To complete the copy $F''$, we must select $6 - v(\tilde{F})$ additional vertices and $6 - e(\tilde{F})$ additional edges. Note that we must have $v(\tilde{F}) > 3$, as otherwise the two vertices in $f'' \cap V(\tilde{F})$ would be contained in two distinct edges in $F''$, a contradiction to the fact that $F$ is linear (i.e., every pair of vertices is contained in at most one edge). Since $F$ is strictly $3$-balanced and $p \ge n^{-1/m_3(F)}$, then, by Claim \ref{claim:strictly-balanced}, we have $n^{v(\tilde{F})} p^{e(\tilde{F})} \ge n^{3 + \lambda} p$ for some $\lambda > 0$. Thus, 
    \[
        n^{6 - v(\tilde{F})} p^{6 - e(\tilde{F})} \le n^{3 - \lambda} p^{5}.
    \]
    Consequently, the contribution of the $i=2$ case to the sum in \eqref{align:sum-of-low-delta} is
    $O\left(\Delta_2(Q) \cdot n^{3 - \lambda} p^{5}\right) = O\left(n^{4 - \lambda} p^{6}\right),$
where we used that $\Delta_2(Q) = O(np)$ which holds by the definition of $\cQ_1$ and $\cQ_2$.

    For $i=3$, there are only $O(1)$ possibilities for the choice of $f'' \in Q$ in the second sum in \eqref{align:sum-of-low-delta}, as the vertex set of $f''$ is completely contained within $V(F')$. There are likewise $O(1)$ ways to choose the intersection $\tilde{F} \coloneqq F' \cap F''$. To complete the copy $F''$, we must select $7 - v(\tilde{F})$ additional vertices and $7 - e(\tilde{F})$ additional edges. Since $F''$ contains $f''$ and at least one other edge from $F'$, the intersection must satisfy $v(\tilde{F}) > 3$. Applying Claim \ref{claim:strictly-balanced} once more, the strictly $3$-balanced nature of $F$ ensures that $n^{v(\tilde{F})} p^{e(\tilde{F})} \ge n^{3 + \lambda} p$ for some $\lambda > 0$. Thus, 
    \[
        n^{7 - v(\tilde{F})} p^{7 - e(\tilde{F})} \le n^{4 - \lambda} p^{6}.
    \]
    Consequently, the contribution of the $i=3$ case to the sum in \eqref{align:sum-of-low-delta} is $O\left(n^{4 - \lambda} p^{6}\right)$.

    Combining the estimates for each case $i \in \{1,2,3\}$ yields the second item of the lemma. \qedhere

    % Fix $F' \in \cFQL[\ext(\Pi)]$. Let $f' \in Q$ satisfying $\{f'\} \cup F' \cong F$. For every $F'' \in \cFQL[\ext(\Pi)]$, there exists $f'' \in Q$ such that $\{f''\} \cup F'' \cong F$. Then, we have
    
    % We have
    % \begin{align}\label{eq:Delta-sum}
    %     \sum_{\substack{F' \neq F'' \in \cFQL[\ext(\Pi)] \\ F' \cap F'' \neq \emptyset}} \Pr(F'' \subseteq \Gnp \mid F' \subseteq \Gnp) = \sum_{j=0}^{\ell}        
    %     \sum_{U \subseteq V(F')} \sum_{\substack{F' \neq F'' \in \cFQL[\ext(\Pi)] \\ F' \cap F'' \neq \emptyset\\ V(F') \cap V(F'') = U}} \Pr(F'' \subseteq \Gnp \mid F' \subseteq \Gnp).
    % \end{align}

    % Let $U \subseteq V(F')$. If $|U| \le 2$, then there are no such copies $F''$ since $F$ is a $3$-uniform hypergraph.

    % For $|U| = 3$, then there is exactly one edge in $F'[U]$. Let us first choose an edge $f \in Q$ such that $f \cup F'' \cong F$. If $|f \cap U| = 0$, then there are no choices of $F''$ since otherwise $F''$ would have two vertex-disjoint edges ($f$ and $U$). If $|f \cap U| = 1$, then there are at most $3 \Delta_1(Q)$ choices of $f$ and then there are at most $n^2$ choices for the last two vertices in $F''$ and we should have at least $5$ more edges, this gives us a factor of order $\Delta_1(Q) n^2 p^5$. If $|f \cap U| = 2$, then there are no such $F''$ since otherwise $F''$ would contain two edges lying on two vertices (the edges $U$ and $f$) but there are no such edges in $F$. We cannot have $f = U$ since in this case $F' \cap F'' = \emptyset$. Thus, we got
    % \[
    %     \sum_{\substack{U \subseteq V(F') \\ |U| = 3}} \sum_{\substack{F' \neq F'' \in \cFQL[\ext(\Pi)] \\ F' \cap F'' \neq \emptyset,\ V(F') \cap V(F'') = U}} \Pr(F'' \subseteq \Gnp \mid F' \subseteq \Gnp) = O\left(\Delta_1(Q) n^2 p^5\right).
    % \]

    % For $|U| = 4$, then there are at most one edge in $F'[U]$. Let us first choose an edge $f \in Q$ such that $f \cup F'' \cong F$. Similarly to before, we cannot have $|f \cap U| = 0$. If $|f \cap U| = 1$, then there are $O(\Delta_1(Q))$ choices of $f$ and then there are at most $n$ choices for the last vertex in $F''$ and we should have at least $5$ more edges, this gives us a factor of order $\Delta_1(Q) n p^5$. If $|f \cap U| = 2$, then there are at most $\Delta_2(Q)$ choices of $f$ and then there are at most $n^2$ choices for the last two vertices in $F''$ and we should have at least $5$ more edges, this gives us a factor of order $\Delta_2(Q) n^2 p^5$. We cannot have $|f \cap U| = 3$ since in this case $F' \cap F'' = \emptyset$. Thus, we got
    % \[
    %     \sum_{\substack{U \subseteq V(F') \\ |U| = 4}} \sum_{\substack{F' \neq F'' \in \cFQL[\ext(\Pi)] \\ F' \cap F'' \neq \emptyset,\ V(F') \cap V(F'') = U}} \Pr(F'' \subseteq \Gnp \mid F' \subseteq \Gnp) = O\left(\Delta_1(Q) n p^5 + \Delta_2(Q) n^2 p^5\right).
    % \]

    % For $|U| = 5$, then there are at most two edges in $F'[U]$. Let us first choose an edge $f \in Q$ such that $f \cup F'' \cong F$. Similarly to before, we cannot $|f \cap U| = 0$. If $|f \cap U| = 1$, then there are $O(\Delta_1(Q))$ choices of $f$ and we should have at least $4$ more edges, this gives us a factor of order $\Delta_1(Q) p^4$. If $|f \cap U| = 2$, then there are $O(\Delta_2(Q))$ choices of $f$ and then there are at most $n$ choices for the last vertex in $F''$ and we should have at least $4$ more edges, this gives us a factor of order $\Delta_2(Q) n p^4$. If $|f \cap U| = 3$, then there are at most $n^2$ choices for the last two vertices in $F''$ and we should have at least $5$ more edges, this gives us a factor of order $n^2 p^5$. Thus, we got
    % \[
    %     \sum_{\substack{U \subseteq V(F') \\ |U| = 5}} \sum_{\substack{F' \neq F'' \in \cFQL[\ext(\Pi)] \\ F' \cap F'' \neq \emptyset,\ V(F') \cap V(F'') = U}} \Pr(F'' \subseteq \Gnp \mid F' \subseteq \Gnp) = O\left(\Delta_1(Q) p^4 + \Delta_2(Q) n p^4 + n^2 p^5\right).
    % \]

    % For $|U| = 6$, then there are at most four edges in $F'[U]$. Let us first choose an edge $f \in Q$ such that $f \cup F'' \cong F$. We cannot $|f \cap U| \le 1$ in this case. If $|f \cap U| = 2$, then there are $O(\Delta_2(Q))$ choices of $f$ and we should have at least $2$ more edges, this gives us a factor of order $\Delta_2(Q) p^2$. If $|f \cap U| = 3$, then there are at most $n$ choices for the last vertex in $F''$ and we should have at least $3$ more edges, this gives us a factor of order $n p^3$. Thus, we got
    % \[
    %     \sum_{\substack{U \subseteq V(F') \\ |U| = 6}} \sum_{\substack{F' \neq F'' \in \cFQL[\ext(\Pi)] \\ F' \cap F'' \neq \emptyset,\ V(F') \cap V(F'') = U}} \Pr(F'' \subseteq \Gnp \mid F' \subseteq \Gnp) = O\left(\Delta_2(Q) p^2 + n p^3\right).
    % \]

    % For $|U| = 7$, since $F' \neq F''$, we have $O(p)$ choices for $F''$. Hence, plugging the above in \eqref{eq:Delta-sum} and using that $np \ge 1$, $\Delta_1(Q) \ge \Delta_2(Q)$, and $\Delta_1(Q) \ge 1$, we get
    % \[
    %     \sum_{\substack{F' \neq F'' \in \cFQL[\ext(\Pi)] \\ F' \cap F'' \neq \emptyset}} \Pr(F'' \subseteq \Gnp \mid F' \subseteq \Gnp) = O\left(\Delta_1(Q) n^2 p^5 + n^2 p^5 + n p^3 + p\right).
    % \]
    % Since $\Delta_1(Q) \le \tilde{\epsilon} \frac{n^2 p}{ \log n}$ and $p \ge n^{-2/3}$, there exists a constant $C_{Low} > 0$ (which does not depend on anything) such that
    % \[
    %     \sum_{\substack{F' \neq F'' \in \cFQL[\ext(\Pi)] \\ F' \cap F'' \neq \emptyset}} \Pr(F'' \subseteq \Gnp \mid F' \subseteq \Gnp) \le \frac{C_{Low} \cdot \tilde{\epsilon} \cdot n^4 p^6}{\log n},
    % \]
    % finishing the proof of the lemma.
\end{proof}

\end{document}
